\subsection{Network assembly and rhythm emergence}

Progressive integration of synaptic connections reveals which interactions are necessary and sufficient for rhythm generation (Fig.~\ref{fig:assembly}).
When uncoupled, Si3 neurons fire tonically at approximately 50~Hz (modeling Si1 command drive) while Si2 neurons remain quiescent.
Activating fast Si3$\to$Si2 excitation recruits Si2 into tonic firing; adding slow Si2$\to$Si3 inhibition produces spike-frequency modulation within the E/I loops but, consistent with experimental observations \cite{SakuraiKatz2016}, these loops alone do not generate autonomous bursting.

Burst alternation emerges only when reciprocal inhibition between Si3 neurons is added, organizing Si3 activity into 2--3 second bursts separated by hyperpolarized interburst intervals.
Si2 neurons exhibit corresponding burst episodes driven by their Si3 partners.
The Si3 half-center oscillator, combined with the E/I modules, thus provides the phase-switching mechanism that converts continuous E/I dynamics into discrete alternating bursts.

The complete network, incorporating Si2 reciprocal inhibition and electrical coupling, generates stable antiphase alternation with 3--4 second bursts matching experimental recordings.
Importantly, the stable rhythm is robust to the order in which connections are activated, confirming that network topology rather than assembly history determines the final dynamics.
This distributed architecture---where no single component is solely responsible for rhythm generation---predicts resilience to transient disruptions, which we test against experimental recordings.
